\documentclass[11pt]{article}

\usepackage[T1]{fontenc}
\usepackage{amsmath,amssymb,amsthm}
\usepackage[margin=1in]{geometry}
\usepackage[colorlinks=true,linkcolor=blue,urlcolor=blue]{hyperref}

\newtheorem{theorem}{Theorem}
\newtheorem{lemma}[theorem]{Lemma}
\newtheorem{corollary}[theorem]{Corollary}
\newtheorem{proposition}[theorem]{Proposition}
\newtheorem{conjecture}[theorem]{Conjecture}
\theoremstyle{remark}
\newtheorem{remark}[theorem]{Remark}

\newcommand{\R}{\mathbb R}
\newcommand{\norm}[1]{\left\|#1\right\|}
\newcommand{\ip}[2]{\left\langle#1,#2\right\rangle}
\newcommand{\weakto}{\rightharpoonup}
\newcommand{\esssup}{\operatorname*{ess\,sup}}
\setlength{\emergencystretch}{2em}

\title{The Full q4 Defect Aligns the Oseen Entrance\\
and Reaches the Nonlinear Clock}
\author{John Lawson and Codex}
\date{25 July 2026}

\begin{document}
\maketitle

\begin{abstract}
This round resolves the phase-sensitive sub-clock corridor left by the
preceding cross-frequency theorem, inside the conditional
energy-efficient exact q4 branch.  Uniform spatial tightness permits the
terminal Oseen data to exhaust the entire downward kinetic-energy jump
\(d\), not merely one ultraviolet portion.  The resulting entrance
field \(a\) satisfies
\[
 \ip{u(t)}{a(t)}=d,\qquad
 \norm{a(t)}_2^2\leq d,\qquad
 a(t)\weakto0
 \quad(t\uparrow T).
\]
At the selected record endpoints,
\[
 a(s_j)-\bigl(u(s_j)-u_*\bigr)\longrightarrow0
 \quad\hbox{strongly in }L^2.
\]
Consequently \(u-a\to u_*\) strongly, while both the adjoint energy
measure and the primal--adjoint cross measure equal the positive
terminal kinetic-energy defect measure.

Uniformly over every Fourier cutoff, the cross pairing is asymptotic to
the positive spectral energy of \(a\), and at records to that of
\(u(s_j)-u_*\).  The terminal weak-\(L^3\) low-pass clock then forces
the median cross frequency to
\[
 \kappa_j\gtrsim
 \Lambda_j=m_j^{7/3}j^{2/3}.
\]
Thus the previously open sub-clock corridor is empty for this canonical
full-defect entrance.  A clock-scale packet still has geometrically
summable unweighted residence, so q4 and the Clay problem remain open.
Every theorem here is a conditional repository result pending external
review.
\end{abstract}

\section{Epistemic status and hypotheses}

No assertion below derives the exact q4 branch from arbitrary
Clay-admissible data.  Let \(u\) be a smooth finite-energy
Navier--Stokes solution on \(\R^3\times(0,T)\), continued weakly
continuously in \(L^2\) through a putative first singular time \(T\).
Write
\begin{equation}
\label{eq:defect}
 E_-:=\lim_{t\uparrow T}\norm{u(t)}_2^2,\qquad
 u_*:=u(T),\qquad
 d:=E_--\norm{u_*}_2^2>0.
\end{equation}
The conditional q4 chain supplies selected record endpoints
\(s_j\uparrow T\), record amplitudes \(m_j\to\infty\), uniform spatial
tightness, and
\begin{equation}
\label{eq:q4tails}
 \int_{s_j}^T
 \norm{u(t)}_{L^{3,\infty}}^2\,dt
 \lesssim \frac{m_j^{-7/2}}{j},
 \qquad
 T-s_j\lesssim\frac{m_j^{-11/2}}{j}.
\end{equation}
Put
\begin{equation}
\label{eq:wj}
 w_j:=u(s_j)-u_*.
\end{equation}
Weak continuity and the energy limit give
\begin{equation}
\label{eq:wlimits}
 w_j\weakto0,\qquad
 \norm{w_j}_2^2\longrightarrow d.
\end{equation}

\section{A finite-band detector for the entire defect}

\begin{lemma}[Fixed low-frequency compactness]
\label{lem:fixed-low}
For every fixed smooth compactly supported Fourier low pass
\(\Pi_{\leq L}\),
\[
 \norm{\Pi_{\leq L}w_j}_2\longrightarrow0.
\]
\end{lemma}

\begin{proof}
Uniform spatial tightness of \(u(s_j)\), together with
\(u_*\in L^2\), makes \(w_j\) uniformly tight.  The fields
\(\Pi_{\leq L}w_j\) are uniformly bounded in \(H^1\).

They remain uniformly tight after low-pass convolution.  Indeed, split
\(w_j\) into its part inside a large ball and its uniformly small
exterior part.  Young's inequality controls the latter, while the
Schwartz tail of the convolution kernel controls the output of the
interior part outside a still larger ball.  Rellich compactness on
bounded balls and this tail estimate make
\(\{\Pi_{\leq L}w_j\}\) precompact in global \(L^2\).  Its only weak
limit is zero by \eqref{eq:wlimits}; hence the convergence is strong.
\end{proof}

\begin{lemma}[Full-defect detector]
\label{lem:detector}
There are
\[
 1<L_j<\Omega_j<\infty,\qquad L_j\longrightarrow\infty,
\]
and solenoidal finite-band fields
\[
 g_j:=\Pi_{L_j<|\xi|\leq\Omega_j}w_j
\]
such that
\begin{equation}
\label{eq:detector}
 \norm{g_j-w_j}_2\longrightarrow0,\qquad
 \norm{g_j}_2^2\longrightarrow d,\qquad
 \ip{u(s_j)}{g_j}\longrightarrow d.
\end{equation}
\end{lemma}

\begin{proof}
Apply Lemma~\ref{lem:fixed-low} successively at integer cutoffs and
choose a staircase \(L_j\uparrow\infty\) slowly enough that
\[
 \norm{\Pi_{\leq L_j}w_j}_2\longrightarrow0.
\]
For each \(j\), choose \(\Omega_j>L_j\) so that
\(\norm{\Pi_{>\Omega_j}w_j}_2\leq j^{-1}\).
Orthogonality of the three Fourier regions gives
\(\norm{g_j-w_j}_2\to0\).  Therefore
\(g_j\weakto0\) and \(\norm{g_j}_2^2\to d\).
Finally,
\[
\begin{aligned}
 \ip{u(s_j)}{g_j}
 &=
 \ip{u_*+w_j}{g_j}\\
 &=
 \ip{u_*}{g_j}
 +\norm{w_j}_2^2
 +\ip{w_j}{g_j-w_j}
 \longrightarrow d.
\end{aligned}
\]
Smooth annular cutoffs can replace the sharp projections after an
arbitrarily small \(L^2\) approximation.
\end{proof}

\section{The aligned Oseen entrance}

For each \(j\), solve the projected backward Oseen problem
\begin{equation}
\label{eq:adjoint-j}
 \partial_ta_j+\nu\Delta a_j
 +\mathbb P((u\cdot\nabla)a_j)=0,\qquad
 \nabla\cdot a_j=0,\qquad
 a_j(s_j)=g_j.
\end{equation}
The exact identities are
\begin{equation}
\label{eq:aj-energy}
 \norm{a_j(t)}_2^2
 +2\nu\int_t^{s_j}\norm{\nabla a_j(r)}_2^2\,dr
 =\norm{g_j}_2^2
\end{equation}
and
\begin{equation}
\label{eq:aj-pair}
 \ip{u(t)}{a_j(t)}
 =\ip{u(s_j)}{g_j}.
\end{equation}
Compactness on every strict preterminal interval gives, after
extraction,
\[
 a\in L^\infty(0,T;L^2_\sigma)
 \cap L^2(0,T;\dot H^1_\sigma),
\]
solving \eqref{eq:adjoint-j} on \((0,T)\), such that
\begin{equation}
\label{eq:a-pair-trace}
 \ip{u(t)}{a(t)}=d,\qquad
 a(t)\weakto0\quad(t\uparrow T).
\end{equation}

\begin{theorem}[Full-defect terminal alignment]
\label{thm:alignment}
The entrance field satisfies
\begin{equation}
\label{eq:a-upper}
 \norm{a(t)}_2^2\leq d\qquad(t<T),
\end{equation}
\begin{equation}
\label{eq:a-terminal-norm}
 \lim_{t\uparrow T}\norm{a(t)}_2^2=d,
\end{equation}
and
\begin{equation}
\label{eq:record-align}
 \norm{a(s_j)-w_j}_2\longrightarrow0.
\end{equation}
Moreover, if \(b:=u-a\), then
\begin{equation}
\label{eq:b-strong}
 b(t)\longrightarrow u_*
 \quad\hbox{strongly in }L^2
 \quad(t\uparrow T).
\end{equation}
\end{theorem}

\begin{proof}
Weak lower semicontinuity in the compact-time limit and
\eqref{eq:aj-energy} imply
\[
 \norm{a(t)}_2^2
 \leq\liminf_j\norm{a_j(t)}_2^2
 \leq\lim_j\norm{g_j}_2^2=d,
\]
which proves \eqref{eq:a-upper}.

At \(t=s_j\), equations \eqref{eq:a-pair-trace} and
\eqref{eq:wlimits} give
\[
\begin{aligned}
 \ip{w_j}{a(s_j)}
 &=
 \ip{u(s_j)}{a(s_j)}
 -\ip{u_*}{a(s_j)}\\
 &=d-o(1).
\end{aligned}
\]
Since \(\norm{w_j}_2\to\sqrt d\), Cauchy--Schwarz and
\eqref{eq:a-upper} force
\[
 \norm{a(s_j)}_2\longrightarrow\sqrt d.
\]
Consequently
\[
\begin{aligned}
 \norm{a(s_j)-w_j}_2^2
 &=
 \norm{a(s_j)}_2^2+\norm{w_j}_2^2
 -2\ip{a(s_j)}{w_j}\\
 &\longrightarrow d+d-2d=0.
\end{aligned}
\]

The adjoint energy identity makes \(\norm{a(t)}_2\) nondecreasing as
\(t\uparrow T\).  Its limit is at most \(\sqrt d\) by
\eqref{eq:a-upper} and attains \(\sqrt d\) along \(s_j\); this proves
\eqref{eq:a-terminal-norm}.

Finally \(b=u-a\weakto u_*\), while
\[
\begin{aligned}
 \norm{b(t)}_2^2
 &=
 \norm{u(t)}_2^2+\norm{a(t)}_2^2
 -2\ip{u(t)}{a(t)}\\
 &\longrightarrow E_-+d-2d
 =\norm{u_*}_2^2.
\end{aligned}
\]
Weak convergence and convergence of norms prove \eqref{eq:b-strong}.
\end{proof}

\begin{corollary}[Exact terminal splitting]
\label{cor:splitting}
For every \(t<T\),
\begin{equation}
\label{eq:a-tail-energy}
 \norm{a(t)}_2^2
 +2\nu\int_t^T\norm{\nabla a(r)}_2^2\,dr=d
\end{equation}
and
\begin{equation}
\label{eq:ab-tail}
 \ip{a(t)}{b(t)}
 =d-\norm{a(t)}_2^2
 =2\nu\int_t^T\norm{\nabla a(r)}_2^2\,dr.
\end{equation}
\end{corollary}

\begin{proof}
Use the adjoint energy equality between \(t\) and \(r<T\), then let
\(r\uparrow T\) and apply \eqref{eq:a-terminal-norm}.  The second
identity follows from \(b=u-a\) and \(\ip{u}{a}=d\).
\end{proof}

\section{Equality of all terminal defect measures}

Let
\[
 \mathcal E
 :=\operatorname*{w^*\!-\!lim}_j
 |u(s_j)|^2\,dx,\qquad
 \vartheta:=\mathcal E-|u_*|^2\,dx.
\]
The input terminal-dimension theorem gives
\[
 \vartheta\geq0,\qquad
 \vartheta(\R^3)=d,\qquad
 \operatorname{supp}\vartheta\subset\sigma,
\]
where \(\mathcal H^1(\sigma)=0\).

\begin{theorem}[Positive common defect measure]
\label{thm:measures}
Along the selected record sequence,
\begin{equation}
\label{eq:measure-equality}
 |a(s_j)|^2\,dx
 \overset{*}{\rightharpoonup}\vartheta,\qquad
 u(s_j)\cdot a(s_j)\,dx
 \overset{*}{\rightharpoonup}\vartheta.
\end{equation}
Thus the adjoint energy defect and the cross defect are exactly the
positive primal kinetic-energy defect.
\end{theorem}

\begin{proof}
Equation \eqref{eq:record-align} gives
\[
 \bigl\||a(s_j)|^2-|w_j|^2\bigr\|_1
 \leq
 \bigl(\norm{a(s_j)}_2+\norm{w_j}_2\bigr)
 \norm{a(s_j)-w_j}_2
 \longrightarrow0.
\]
Expanding \(|w_j|^2=|u(s_j)-u_*|^2\), using weak convergence against
compactly supported multiples of \(u_*\), and using spatial tightness
shows that \(|w_j|^2dx\overset{*}{\rightharpoonup}\vartheta\).
This proves the first limit.

For the second,
\[
 u(s_j)\cdot a(s_j)-|w_j|^2
 =
 w_j\cdot(a(s_j)-w_j)+u_*\cdot a(s_j).
\]
The first term tends to zero in \(L^1\).  The second tends weakly to
zero as a measure because \(a(s_j)\weakto0\).  Its tails are uniformly
small by Cauchy--Schwarz and \(u_*\in L^2\).  This proves the second
limit.
\end{proof}

\section{Uniform positivity of the cross spectrum}

Choose a real radial
\(\psi\in C_c^\infty(\R^3)\), nonincreasing in \(|\xi|\), with
\[
 0\leq\psi\leq1,\qquad
 \psi=1\quad(|\xi|\leq1),\qquad
 \psi=0\quad(|\xi|\geq2),
\]
and set \(B_N:=\psi(D/N)\).  Define
\begin{equation}
\label{eq:KN}
 K_N(t):=\ip{u(t)}{B_Na(t)}.
\end{equation}

\begin{lemma}[Compact cutoff orbit]
\label{lem:cutoff-orbit}
For every \(f\in L^2\), the set
\[
 \{B_Nf:0<N<\infty\}\cup\{0,f\}
\]
is norm-compact in \(L^2\).
\end{lemma}

\begin{proof}
The map \(N\mapsto B_Nf\) is norm-continuous by dominated convergence,
tends to zero as \(N\downarrow0\), and tends to \(f\) as
\(N\uparrow\infty\).  It therefore extends continuously to a compact
parameter interval.
\end{proof}

\begin{theorem}[Asymptotically positive cross spectrum]
\label{thm:positive-spectrum}
Uniformly over all \(N>0\),
\begin{equation}
\label{eq:a-spectrum}
 K_N(t)=\ip{a(t)}{B_Na(t)}+o_{t\uparrow T}(1),
\end{equation}
and, along the records,
\begin{equation}
\label{eq:w-spectrum}
 K_N(s_j)=\ip{w_j}{B_Nw_j}+o_{j\to\infty}(1).
\end{equation}
\end{theorem}

\begin{proof}
Since \(u=a+b\),
\[
 K_N-\ip{a}{B_Na}=\ip{b}{B_Na}.
\]
The part with \(b-u_*\) is uniformly bounded by
\(\norm{b-u_*}_2\norm a_2\), which tends to zero by
\eqref{eq:b-strong}.  The remaining part is
\(\ip{B_Nu_*}{a}\).  Weak convergence \(a(t)\weakto0\) is uniform on
the compact set from Lemma~\ref{lem:cutoff-orbit}, proving
\eqref{eq:a-spectrum}.

At \(s_j\),
\[
\begin{aligned}
 K_N(s_j)-\ip{w_j}{B_Nw_j}
 ={}&
 \ip{u(s_j)}{B_N(a(s_j)-w_j)}\\
 &+\ip{B_Nu_*}{w_j}.
\end{aligned}
\]
The first term tends uniformly to zero by
\eqref{eq:record-align}; the second does so because
\(w_j\weakto0\) uniformly on the same compact cutoff orbit.
\end{proof}

\begin{remark}
The functions \(N\mapsto\ip f{B_Nf}\) are nonnegative and
nondecreasing.  Thus Theorem~\ref{thm:positive-spectrum} removes, up to
a uniform \(o(1)\), the sign, phase, and nonmonotonicity ambiguity that
created the earlier corridor.
\end{remark}

\section{The terminal low-pass clock}

Put
\[
 M(t):=\norm{u(t)}_{L^{3,\infty}},\qquad
 B_{2,j}:=\int_{s_j}^TM(t)^2\,dt,\qquad
 \tau_j:=T-s_j,
\]
and \(U:=\sup_{t<T}\norm{u(t)}_2\).

\begin{lemma}[Terminal weak-\(L^3\) low-pass estimate]
\label{lem:terminal-clock}
For every \(N>0\),
\begin{equation}
\label{eq:terminal-clock}
 \norm{B_Nw_j}_2
 \lesssim
 N^{3/2}B_{2,j}
 +\nu U N^2\tau_j.
\end{equation}
\end{lemma}

\begin{proof}
Apply \(B_N\) to the projected Navier--Stokes equation.  Lorentz
duality and Bernstein give
\[
 \norm{B_N\mathbb P\nabla\cdot(u\otimes u)}_2
 \lesssim
 N^{3/2}\norm{u\otimes u}_{L^{3/2,\infty}}
 \lesssim N^{3/2}M(t)^2.
\]
The viscous term satisfies
\[
 \norm{B_N\Delta u(t)}_2\lesssim N^2U.
\]
Integrate from \(s_j\) to \(r<T\).  For fixed \(N\), the same
integrable estimate makes \(B_Nu(r)\) strongly Cauchy as
\(r\uparrow T\).  Its weak limit identifies the strong limit as
\(B_Nu_*\).  Letting \(r\uparrow T\) proves
\eqref{eq:terminal-clock}.
\end{proof}

Define the first dyadic half-pairing frequency
\begin{equation}
\label{eq:kappa}
 \kappa_j
 :=
 \min\{2^k:K_{2^k}(s_j)\geq d/2\}.
\end{equation}

\begin{theorem}[Sub-clock corridor collapse]
\label{thm:corridor-collapse}
For all sufficiently large \(j\),
\begin{equation}
\label{eq:clock-frequency}
 \boxed{
 \kappa_j\gtrsim
 \Lambda_j:=m_j^{7/3}j^{2/3}.
 }
\end{equation}
\end{theorem}

\begin{proof}
The uniform alignment \eqref{eq:w-spectrum} gives
\[
 \ip{w_j}{B_{\kappa_j}w_j}\geq d/3
\]
for large \(j\).  Since \(B_{2\kappa_j}=1\) on the Fourier support of
the symbol of \(B_{\kappa_j}\),
\[
\begin{aligned}
 \frac d3
 &\leq
 \ip{B_{2\kappa_j}w_j}{B_{\kappa_j}w_j}\\
 &\leq
 \norm{B_{2\kappa_j}w_j}_2\norm{w_j}_2.
\end{aligned}
\]
The \(w_j\) norms tend to \(\sqrt d\).  Applying
Lemma~\ref{lem:terminal-clock} at \(2\kappa_j\) yields
\[
 1\lesssim
 \kappa_j^{3/2}B_{2,j}
 +\nu U\kappa_j^2\tau_j.
\]
If the first term pays, \eqref{eq:q4tails} gives
\[
 \kappa_j\gtrsim B_{2,j}^{-2/3}
 \gtrsim m_j^{7/3}j^{2/3}.
\]
If the second pays, then
\[
 \kappa_j\gtrsim\tau_j^{-1/2}
 \gtrsim m_j^{11/4}j^{1/2}.
\]
The latter exceeds \(\Lambda_j\), because
\[
 \frac{m_j^{11/4}j^{1/2}}
 {m_j^{7/3}j^{2/3}}
 =
 m_j^{5/12}j^{-1/6}\longrightarrow\infty.
\]
Both cases prove \eqref{eq:clock-frequency}.
\end{proof}

\begin{remark}
The proof uses more than the absolute commutator inequalities: it uses
full capture of the kinetic defect, strong primal--adjoint alignment,
and the terminal Navier--Stokes low-pass evolution.  The scalar
sub-clock front from the preceding round lacks precisely these three
facts.
\end{remark}

\section{Individual enstrophy floors}

Use reverse terminal time \(\tau=T-t\), and put
\[
 M(\tau):=\norm{u(T-\tau)}_{L^{3,\infty}},\quad
 X(\tau):=\norm{\nabla u(T-\tau)}_2,\quad
 Y(\tau):=\norm{\nabla a(T-\tau)}_2,
\]
\[
 B_1(\tau):=\int_0^\tau M(r)\,dr.
\]
The preceding coarse commutator theorem gives
\begin{equation}
\label{eq:coarse-front}
 \kappa(\tau)^2B_1(\tau)\gtrsim1.
\end{equation}

\begin{proposition}[Individual moving-front cost]
\label{prop:individual}
For all sufficiently small \(\tau\),
\begin{equation}
\label{eq:individual}
 X(\tau)^2\gtrsim\kappa(\tau)^2
 \gtrsim B_1(\tau)^{-1},
 \qquad
 Y(\tau)^2\gtrsim\kappa(\tau)^2
 \gtrsim B_1(\tau)^{-1}.
\end{equation}
If
\(\ell(\tau):=\log(e+\tau_0/\tau)\), then
\begin{equation}
\label{eq:individual-tails}
 \int_0^\tau X(r)^2\,dr
 \gtrsim\tau^{2/11}\ell(\tau)^{2/11},
 \qquad
 \int_0^\tau Y(r)^2\,dr
 \gtrsim\tau^{2/11}\ell(\tau)^{2/11}.
\end{equation}
Moreover,
\begin{equation}
\label{eq:weighted-both}
 \int_0^{\tau_0}M(\tau)X(\tau)^2\,d\tau=\infty,
 \qquad
 \int_0^{\tau_0}M(\tau)Y(\tau)^2\,d\tau=\infty.
\end{equation}
\end{proposition}

\begin{proof}
At the dyadic predecessor \(N=\kappa(\tau)/2\), minimality and the full
pairing \(d\) give
\[
 \frac d2<\ip{u}{(I-B_N)a}.
\]
The high-pass estimate
\[
 \norm{(I-B_N)f}_2\lesssim N^{-1}\norm{\nabla f}_2
\]
can be applied on either side of the pairing.  Together with
\(\norm a_2\leq\sqrt d\) and the uniform \(L^2\) bound for \(u\), it
gives separately
\[
 X(\tau)\gtrsim\kappa(\tau),\qquad
 Y(\tau)\gtrsim\kappa(\tau).
\]
Combine this with \eqref{eq:coarse-front}.

The q4 clock gives
\[
 B_1(\tau)
 \lesssim
 \tau^{9/11}\ell(\tau)^{-2/11}.
\]
Integrating \(B_1^{-1}\) proves
\eqref{eq:individual-tails}.  Finally \(B_1'=M\) almost everywhere, so
\[
 \int_\epsilon^{\tau_0}\frac{M(\tau)}{B_1(\tau)}\,d\tau
 =
 \log B_1(\tau_0)-\log B_1(\epsilon)
 \longrightarrow\infty
\]
as \(\epsilon\downarrow0\).  Equation \eqref{eq:individual} proves
\eqref{eq:weighted-both}.
\end{proof}

\begin{corollary}
At the selected q4 records,
\[
 \norm{\nabla u(s_j)}_2^2\gtrsim\Lambda_j^2,\qquad
 \norm{\nabla a(s_j)}_2^2\gtrsim\Lambda_j^2.
\]
\end{corollary}

\section{Sharp clock-scale survivor}

Theorem~\ref{thm:corridor-collapse} is not q4 closure.  Take the exact
representative powers
\[
 m_j:=2^{2j},\qquad
 \delta_j:=\frac{2^{-11j}}j,\qquad
 \Lambda_j:=2^{14j/3}j^{2/3}.
\]
Then
\begin{equation}
\label{eq:unweighted-ledger}
 \delta_j\Lambda_j^2
 =2^{-5j/3}j^{1/3},
 \qquad
 \sum_j\delta_j\Lambda_j^2<\infty,
\end{equation}
whereas
\begin{equation}
\label{eq:weighted-ledger}
 m_j\delta_j\Lambda_j^2
 =2^{j/3}j^{1/3},
 \qquad
 \sum_jm_j\delta_j\Lambda_j^2=\infty.
\end{equation}
Thus a packet moving at the nonlinear clock can retain finite primal
and adjoint unweighted dissipation while satisfying the forced
amplitude-weighted divergence.  This is a temporal-frequency scaling
ledger, not a Navier--Stokes construction.

\section{Round audit}

\subsection*{Formal conditional results}

\begin{enumerate}
\item Finite annular terminal data can exhaust the entire kinetic defect
      \(d\).
\item The resulting Oseen entrance obeys
      \(\norm a_2^2\leq d\), has terminal norm exactly \(\sqrt d\), and
      aligns strongly with \(u(s_j)-u_*\).
\item The complementary field \(u-a\) reaches \(u_*\) strongly.
\item The adjoint energy measure and cross measure both equal the
      positive kinetic defect \(\vartheta\).
\item The cross spectrum is uniformly asymptotic to a positive energy
      spectrum.
\item The terminal low-pass clock forces
      \(\kappa_j\gtrsim\Lambda_j\), closing the sub-clock corridor for
      the full-defect entrance.
\item The individual enstrophy tails satisfy the stronger
      \(\tau^{2/11}\ell(\tau)^{2/11}\) floors, and both weighted actions
      diverge.
\end{enumerate}

\subsection*{Closed proposal}

Arbitrary phase or sign is not the q4 survivor.  The canonical entrance
is positively aligned with the whole primal kinetic defect, and its
cross front reaches the direct nonlinear clock.  Summing the resulting
pointwise clock floors is nevertheless insufficient because
\eqref{eq:unweighted-ledger} is geometrically summable.

\subsection*{Open obligations}

\begin{enumerate}
\item Turn the clock-scale defect quantile into a non-reusable
      residence, flux, or projective-motion charge.
\item Exploit the positive identity between the adjoint, cross, and
      primal defect measures in the local current or capacity equation.
\item Use the strong component \(u-a\to u_*\) to obtain a coercive
      relative-energy or projected-pressure identity.
\item Prove strong left \(L^2\) continuity, energy equality, finite
      weighted adjoint action, or another established closing criterion.
\item Treat slower clocks, divergent normalised energy, and the other
      Clay alternatives separately.
\end{enumerate}

\begin{conjecture}[No clock-scale full-defect entrance]
No smooth finite-energy Navier--Stokes trajectory approaching its first
singular time possesses a full-defect projected-Oseen entrance \(a\)
such that
\[
 a(t)\weakto0,\qquad
 u(t)-a(t)\to u_*\ \hbox{strongly},\qquad
 \kappa_j\gtrsim\Lambda_j,
\]
while the common positive defect measure is supported on an
\(\mathcal H^1\)-null terminal singular slice.
\end{conjecture}

The conjecture is not proved.  This round closes only the sub-clock
corridor inside the conditional energy-efficient exact q4 branch.  No
alternative A--D of the Clay problem is proved.

\end{document}
